% ── Section 6: Discussion ───────────────────────────────────────────────────
\section{Discussion}
\label{sec:discussion}

The results highlight the varying adaptive capacities of reinforcement learning procurement policies operating under non-stationary market regimes. Across the evaluated configurations, the MLP showed the least stable performance, whereas the LSTM provided more conservative behaviour. The GTrXL configurations achieved higher mean savings in the Walk-Forward evaluation, particularly when combined with the NLP features.

The Multi-Layer Perceptron (MLP) demonstrated unstable performance during the static 2025 evaluation, suggesting that its lack of temporal memory restricts adaptation to distributional shifts. Conversely, the RecurrentPPO (LSTM) agent converged to a single-contract policy under static evaluation, prioritising the D2 contract. Under the Walk-Forward evaluation without NLP data, the LSTM Quantitative Commodities Baseline model yielded mean losses ($-2.00\%$), indicating that sequence memory alone was insufficient to maintain positive savings given only quantitative features.

\subsection{Characterisation of Low-Entropy Policy Convergence}

In the evaluated static runs, the recurrent policy converged to a low-entropy behaviour dominated by D2 selections. This concentration reduced exposure to longer-horizon contracts but also limited the range of procurement actions available to the policy. This behaviour is compatible with the combined effect of delayed reward realization at delivery, non-overlapping contract penalties, and exponential temporal discounting ($\gamma=0.99$), although our experiments do not isolate the individual contribution of each mechanism.

The Walk-Forward findings indicate that this low-entropy convergence is mitigated when the state representation includes multimodal information under continuous retraining. The Hybrid representation (\texttt{Dataset\_M8\_Physical\_NLP}), combining physical fundamentals with targeted macroeconomic sentiment, yields a more balanced action profile (D2: 51.6\%, D1: 17.9\%, SPOT: 16.6\%, D3: 13.9\%) while maintaining positive simulated savings during the 2026 shock (+6.32\%).

A central theoretical consideration is the functional role of zero-shot LLM news sentiment as an alternative representation. Under stationary conditions, quantitative commodity prices and historical generation curves provide direct econometric signals for forward pricing. However, when wholesale market dislocations occur---such as the renewable price cannibalisation observed in Iberia in 2026---historical empirical correlations between international fossil fuel prices (Brent, TTF, EU ETS) and local electricity clearing prices may decouple. The sentiment variables offered an alternative representation of regulatory, geopolitical, and generation-related information. Their inclusion coincided with positive savings (+6.32\%) during the 2026 evaluation period, although the present experiments do not establish whether this improvement resulted primarily from semantic narrative signals, feature dimensionality, noise filtering, or interactions with the policy architecture. Establishing a formal causal regime-filtering mechanism would require an independently held-out multi-year period and explicit econometric time-varying-parameter tests.

This progression is reflected in the Phase 4 results. The sentiment variables added information that was not strongly collinear with several physical features in the descriptive PCA and were associated with higher procurement savings in the evaluated configurations. Augmenting the state representation with lagged news-derived sentiment on LNG availability, regulatory developments, and geopolitical risk yielded $+9.36\%$ mean savings ($\sigma=1.37\%$) for the LSTM Hybrid model, compared to the quantitative commodity baseline ($-2.00\%$, $\sigma=1.56\%$). This difference is consistent with an incremental contribution from qualitative news sentiment relative to numerical commodity futures during the evaluated shift.

In the evaluated setting, the GTrXL configuration achieved higher mean savings than the LSTM configuration. The Transformer Quantitative Commodities model reached $+9.97\%$ ($\sigma=1.32\%$). The Hybrid Transformer model yielded $+10.85\%$ mean savings, accompanied by a reduction in across-seed standard deviation to $\sigma=0.63\%$. Across-seed standard deviation ($\sigma_{\text{seeds}}$) quantifies algorithmic optimization stability under stochastic neural initialization rather than downside financial risk (e.g., CVaR or maximum drawdown). Future industrial extensions should incorporate explicit risk-averse utility functions (such as distributional RL or spectral risk penalties) to directly penalize tail loss events during market spikes.

\subsection{Economic Evaluation: Dynamic Multi-Horizon DRL vs.\ Partial Hedging and Static Heuristics}
\label{subsec:economic_evaluation}

An essential operational and scientific question is why an industrial consumer or energy procurement manager should deploy a complex reinforcement learning policy rather than executing a simple passive procurement heuristic---such as systematic $D_2$ purchasing or a static partial hedging rule ($\alpha \cdot D_2 + (1-\alpha) \cdot \text{SPOT}$). Across the annual Walk-Forward evaluations (Table~\ref{tab:annual_rl_vs_d2}) and the 2026 isolated shock (Table~\ref{tab:blind_2026_comparison}), the mechanical $100\%$ $D_2$ heuristic achieved higher nominal savings ($+11.18\%$ to $+15.67\%$) than the evaluated RecurrentPPO policy ($+5.45\%$ to $+6.74\%$). During the isolated January--May 2026 evaluation, M8-W1 achieved $6.32\%$ simulated savings, compared with $11.18\%$ for the passive $D_2$ rule and $5.59\%$ for a $50\% D_2 / 50\% \text{SPOT}$ hedge. The policy had lower downside semi-deviation than $D_2$ ($9.15$ versus $13.68$~\texteuro/MWh) and a lower maximum drawdown ($186.45$ versus $225.45$~\texteuro/MWh), but its paired savings differences against $D_2$ and the balanced hedge had moving-block bootstrap confidence intervals that included zero. The results therefore suggest a different savings--risk profile in this episode, rather than statistically established savings superiority. The 2026 window represents a single market-dislocation episode and cannot establish general regime robustness.

This empirical divergence raises fundamental economic considerations regarding the true objective of procurement policies in volatile energy markets. We delineate four critical operational and risk-management considerations regarding dynamic reinforcement learning relative to static heuristics:

\subsubsection{Historical Sample Bias and Forward--Spot Differentials}

The high nominal savings of the mechanical $D_2$ rule in the evaluated sample are consistent with a horizon-dependent forward--spot differential, but the present data do not identify its cause. Risk premia, market-maker quotes, liquidity conditions, delivery-horizon effects, and sample-specific market dynamics may all contribute~\cite{huisman2012electricity, woo2001cross}. In the absence of direct participant, bid--ask, order-book, and renewable-forecast-error data, we do not attribute the observed differential specifically to renewable forecast errors or to a causal market mechanism.

Relying on the persistence of a single market anomaly represents an unhedged speculative bet on market microstructure stationarity. When market rules change, when forecasting tools evolve, or when liquidity migrates, historical forward discounts can invert rapidly. In our empirical sample, identical mechanical purchasing on $D_1$ yielded negative savings ($-1.30\%$), $D_3$ suffered high downside volatility ($\sigma_- = 18.63$~\texteuro/MWh) and substantial drawdown ($548.95$~\texteuro/MWh), and longer-term monthly contracts collapsed to severe losses ($-20.72\%$ for $M_3$, $-31.61\%$ for $M_4$, and $-49.95\%$ in 2023). A static heuristic possesses no mechanism to detect when a discount has evaporated; it continues executing blindly until major drawdown occurs.

\subsubsection{Comparison against the Partial Hedging Continuum \texorpdfstring{($\alpha \cdot D_2 + (1-\alpha) \cdot \text{SPOT}$)}{(D2--SPOT Combinations)}}

Industrial risk managers rarely execute $100\%$ unhedged SPOT or $100\%$ single-instrument forwards; instead, corporate hedging mandates typically prescribe linear partial hedging mixtures $\alpha \in [0, 1]$. As formalized in Panel~A of Table~\ref{tab:blind_2026_comparison}, a static partial hedge linearly interpolates both return and downside risk ($g_t(\alpha) = \alpha \cdot g_t(D_2)$). A balanced $50\%\, D_2 + 50\%\, \text{SPOT}$ policy achieves $+5.59\%$ savings with a maximum drawdown of $112.73$~\texteuro/MWh.

The dynamic DRL agent (M8-W1) achieves a point estimate of $+6.32\%$ savings, compared with $+5.59\%$ for the balanced $50\%$ partial hedge ($+0.73$~pp difference, $d_t = +0.32$~\texteuro/MWh). However, because the paired 95\% moving-block bootstrap confidence interval spans $[-3.10, +5.91]$~pp, this difference does not establish statistically significant savings superiority over the balanced hedge. Instead, the dynamic policy demonstrates an alternative savings--risk profile: rather than maintaining an inflexible static ratio, the DRL agent dynamically modulates its allocation conditional on daily market fundamentals and macroeconomic news sentiment: holding SPOT unhedged when pool prices are projected to collapse (16.6\% of sessions), and dynamically selecting between $D_1$ (17.9\%), $D_2$ (51.6\%), and $D_3$ (13.9\%) depending on where the forward risk premium is concentrated. In the Phase~4 capacity evaluation, the GTrXL Transformer agent achieved $+10.85\%$ savings, matching the nominal return level of mechanical $D_2$ ($+11.18\%$) while maintaining a diversified multi-horizon execution profile.

\subsubsection{Action Diversification across Horizons vs.\ Static Concentration}

A notable operational consideration regarding the mechanical $100\%$ $D_2$ strategy is single-instrument market concentration. On organized derivative exchanges such as OMIP, short-dated daily forward contracts ($D_1, D_2, D_3$) exhibit finite order-book liquidity and depth. For a large industrial enterprise or aggregator consuming, for example, 50 to 200~MW of continuous baseload ($1{,}200$ to $4{,}800$~MWh/day), attempting to purchase $100\%$ of required delivery volume exclusively on the $D_2$ order book every morning would concentrate execution exposure in a single instrument.

In contrast, the dynamic DRL policy demonstrates simulated action diversification across horizons: order flow is distributed across $D_1$ (17.9\%), $D_2$ (51.6\%), $D_3$ (13.9\%), and SPOT daily clearing auctions (16.6\%). We caution that our simulation environment does not model order-book depth, bid--ask spreads, or non-linear market impact; therefore, this multi-horizon distribution demonstrates policy flexibility and diversification within the simulation, rather than empirical proof of reduced market impact in physical execution.

\subsubsection{Regime-Adaptive Risk Management vs.\ Static Fragility}

Ultimately, the primary scientific and economic finding of the proposed framework is not claiming unconditional return dominance over an isolated historical anomaly, but evaluating an \textbf{adaptive multi-horizon risk mitigator}. In energy procurement, downside risk---quantified by Downside Semi-deviation ($\sigma_-$), Conditional Value at Risk ($\text{CVaR}_{0.95}$), and Maximum Drawdown---is often more critical for corporate solvency than nominal return maximization.

During the 2026 wholesale price dislocation, the dynamic M8-W1 policy reduced downside semi-deviation from $13.68$ to $9.15$~\texteuro/MWh (for $D_2$, a $33.1\%$ volatility reduction) and lowered maximum drawdown from $225.45$ to $186.45$~\texteuro/MWh. Compared to alternative forward contracts such as $D_3$ (Max DD $= 548.95$~\texteuro/MWh), the RL policy cut maximum drawdown by nearly two-thirds in this sample. By combining physical generation state variables with forward-looking macroeconomic news sentiment, the DRL agent actively navigates non-stationary market regimes, offering corporate buyers an alternative multi-horizon hedging profile that curtails extreme tail losses when historical market correlations breakdown.
\subsection{Transaction Cost Sensitivity Analysis}
\label{subsec:trans_costs}

The primary backtest evaluations assume frictionless forward execution ($C_0 = 0.00$~\texteuro/MWh) to isolate the theoretical settlement-implied procurement efficiency of the learned policy. In industrial wholesale electricity trading, continuous hedging via broker-intermediated OTC contracts or organized derivatives exchanges (such as OMIP and OMIClear) incurs clearing fees, exchange registration tariffs, and bid--ask half-spreads. While general power market literature reports that execution tariffs and effective spreads typically range between $0.05$ and $0.25$~\texteuro/MWh~\cite{weron2014electricity}, exact costs depend on consumer volume, exchange membership, and broker bilateral arrangements.

To assess policy sensitivity under market frictions, we conducted an exploratory sensitivity analysis across three illustrative linear transaction-cost stress scenarios applied to all executed forward contracts: a frictionless baseline ($C_0 = 0.00$~\texteuro/MWh), a moderate illustrative friction tier ($C_1 = 0.10$~\texteuro/MWh), and an adverse friction stress scenario ($C_2 = 0.25$~\texteuro/MWh). We emphasize that these are synthetic linear sensitivity tiers rather than costs observed directly from OMIP order books.

The arithmetic of the cost impact is formally derived from the empirical delivery settlement log. Percentage procurement savings relative to spot clearing are defined as $\text{Savings} = (\bar{P}_{\text{SPOT}} - \bar{P}_{\text{effective}}) / \bar{P}_{\text{SPOT}}$. When an execution fee $C$ (\texteuro/MWh) is incurred on forward contracts, effective procurement unit expenditure rises by $f_{\text{fwd}} \cdot C$, where $f_{\text{fwd}}$ represents the empirical fraction of delivered MWh settled under locked forward positions (with the residual $1 - f_{\text{fwd}}$ defaulting to unhedged spot clearing without forward execution fees). From the multi-year Walk-Forward decision and delivery log across matched evaluation cohorts, forward commitments covered exactly $f_{\text{fwd}} = 85.8\%$ of delivery days ($321$ of 374 evaluated days: $53.5\%$ under $D_2$, $17.6\%$ under $D_1$, and $14.7\%$ under $D_3$), while unhedged spot clearing settled the remaining $14.2\%$ ($53$ days), against an empirical mean spot price of $\bar{P}_{\text{SPOT}} = 65.76$~\texteuro/MWh across those 374 evaluated days. In the early-2026 Walk-Forward test blocks (44 evaluated days), forward delivery commitments accounted for $f_{\text{fwd}} = 75.0\%$ ($33$ of 44 days under $D_2$) and spot clearing settled $25.0\%$ ($11$ days), against a mean spot price of $\bar{P}_{\text{SPOT}} = 36.41$~\texteuro/MWh across those 44 evaluated days. The resulting percentage reduction in savings is $\Delta \text{Savings} = (f_{\text{fwd}} \cdot C) / \bar{P}_{\text{SPOT}}$:
\begin{itemize}
    \item \textbf{Full Walk-Forward Horizon} (374 days, $\bar{P}_{\text{SPOT}} = 65.76$ \texteuro/MWh, $f_{\text{fwd}} = 85.8\%$): For $C_1 = 0.10$ \texteuro/MWh, the expected friction is $0.0858$ \texteuro/MWh, reducing net savings by $0.0858 / 65.76 \approx 0.13$~pp ($+6.04\% \to +5.91\%$). For $C_2 = 0.25$ \texteuro/MWh, the friction is $0.2146$ \texteuro/MWh, reducing net savings by $0.2146 / 65.76 \approx 0.33$~pp ($+6.04\% \to +5.71\%$).
    \item \textbf{2026 Walk-Forward Test Blocks} (44 days, $\bar{P}_{\text{SPOT}} = 36.41$ \texteuro/MWh, $f_{\text{fwd}} = 75.0\%$): For $C_1 = 0.10$ \texteuro/MWh, the expected friction is $0.0750$ \texteuro/MWh, reducing net savings by $0.0750 / 36.41 \approx 0.21$~pp ($+6.32\% \to +6.11\%$). For $C_2 = 0.25$ \texteuro/MWh, the friction is $0.1875$ \texteuro/MWh, reducing net savings by $0.1875 / 36.41 \approx 0.51$~pp ($+6.32\% \to +5.80\%$).
\end{itemize}

\begin{table*}[htbp]
\centering
\footnotesize
\setlength{\tabcolsep}{3.5pt}
\caption{Sensitivity of Walk-Forward procurement savings to illustrative linear forward execution transaction costs ($C$, \texteuro/MWh) based on empirical delivery settlement logs across matched evaluation horizons.}
\label{tab:trans_cost_sensitivity}
\begin{tabularx}{\linewidth}{@{}l >{\raggedright\arraybackslash}X ccc@{}}
\toprule
\textbf{Cost Tier} & \textbf{Illustrative Cost Description} & \textbf{Full Horizon (374 d)} & \textbf{2026 WF (44 d)} & \textbf{Net} \\
\textbf{($C$, \texteuro/MWh)} & & \textbf{Savings (\%)} & \textbf{Savings (\%)} & \textbf{Viability} \\
\midrule
$C_0 = 0.00$ & Frictionless baseline & $+6.04\%$ & $+6.32\%$ & Baseline \\
$C_1 = 0.10$ & Moderate illustrative friction & $+5.91\%$ & $+6.11\%$ & Positive \\
$C_2 = 0.25$ & Adverse illustrative stress tier & $+5.71\%$ & $+5.80\%$ & Positive \\
\bottomrule
\end{tabularx}
\end{table*}

As shown in Table~\ref{tab:trans_cost_sensitivity}, simulated net procurement savings remain positive across all tested linear cost scenarios. Under the adverse friction assumption ($C_2 = 0.25$~\texteuro/MWh), the Walk-Forward M8 policy retains $+5.71\%$ net savings across the full evaluation horizon (374 days) and $+5.80\%$ during the 2026 Walk-Forward test blocks (44 days). Under the illustrative linear-cost assumptions examined, simulated savings remain positive. These calculations are a sensitivity analysis rather than an estimate of realized execution costs, because the environment does not model order-book depth, bid--ask spreads, or non-linear market impact.

\subsection{Limitations and Threats to Validity}

Despite the empirical performance, several limitations must be acknowledged:

First, holding total monthly PPO training interactions fixed at $K = 15{,}000$ steps across all sliding window lengths reflects an operational computational budget constraint. A short update window ($W=1$ month, $N_W \approx 20$ sessions) receives $\approx 750$ gradient interaction steps per historical observation, a quarterly window ($W=3$ months) receives $\approx 250$ steps/obs, a semi-annual window ($W=6$ months) receives $\approx 125$ steps/obs, whereas an annual window ($W=12$ months, $N_W \approx 240$ sessions) receives $\approx 62.5$ steps per observation. While fixing $K$ represents an operationally viable deployment protocol with deterministic monthly retraining latency, an alternative experimental design would scale $K$ proportionally with window size ($K(W) = 750 \times N_W$, scaling up to $K \approx 180{,}000$ for $W=12$) to hold optimization intensity per observation constant. Future work should evaluate this proportional scaling ablation to formally decouple sample recency from per-sample gradient budget.

Second, the choice of discount factor $\gamma = 0.99$ introduces temporal discounting of distant cash flows ($\gamma^{30} \approx 0.74, \gamma^{90} \approx 0.40, \gamma^{180} \approx 0.16$). In conjunction with single-instrument full-volume baseload commitments and non-overlapping forward masking, geometric discounting may comparatively reduce the learning gradient signal associated with distant delivery horizons ($M_1$--$M_6$), guiding policy optimization toward contracts with immediate delivery proximity ($D_1$--$D_3$). The observed policy concentration on short-horizon products is therefore jointly shaped by the empirical forward--spot risk premia and temporal MDP mechanics.

Third, regarding temporal integrity, the LLM sentiment extraction was executed locally using the Llama-3 model with zero internet access during inference, ensuring data privacy, absence of external API drift, and strict computational repeatability under deterministic decoding. While deterministic decoding under the reported software configuration ($T=0.0$, top-$p=1.0$, strict JSON schema) enforces syntactic and computational repeatability, future work could further substantiate semantic construct validity by benchmarking against an independent human-annotated corpus. Measuring inter-annotator agreement (using weighted Cohen's $\kappa$ for ordinal categories and Intraclass Correlation Coefficients or Spearman rank correlations for continuous indices) across domain-expert annotations would provide additional external validation for specialized energy-market sentiment models. Furthermore, to conclusively demonstrate that the performance gains of the Hybrid representation stem from genuine semantic directionality rather than implicit dimensional regularisation (i.e., expanding the feature vector from 35 to 45 dimensions), a future ablation could substitute the 10 NLP sentiment signals ($\phi_1, \dots, \phi_{10}$) with matched Gaussian noise $\boldsymbol{\epsilon} \sim \mathcal{N}(\boldsymbol{\mu}_{\text{NLP}}, \boldsymbol{\Sigma}_{\text{NLP}})$. Nevertheless, the SHAP feature attribution analysis in Section~\ref{subsec:res_shap} confirms that the policy logits systematically attend to domain-coherent news dimensions (such as regulatory risk and generation mix), corroborating the semantic grounding of the extracted features.

Fourth, Phases 1 and 2 report single-seed point estimates across the 8-dataset ablation space. While computationally necessary to explore the broad factorial design space before committing to the intensive multi-seed retraining protocol, the descriptive interaction observed in Phase~1 remains preliminary until multi-seed replications and dimension-matched noise controls across all factorial cells are executed. Multi-seed validation was rigorously performed for all continuous Walk-Forward and capacity evaluations in Phases 3 and 4 across 5 independent seeds.

Fifth, in alignment with standard machine-learning validation protocols, the 2025 horizon served as the preliminary validation partition to discard unviable heuristics and prune unpromising feature representations. Consequently, the multi-year continuous Walk-Forward evaluation preserves the 2026 horizon (January--May 2026, comprising 102 official OMIP trading sessions [decision epochs] and 46 non-trading dates [3 TARGET closures and 43 weekend days] within 148 calendar delivery days) as a strictly forward-looking, sequentially out-of-sample benchmark (with each monthly decision conditioned only on information available prior to that evaluation block) to test resilience under the wholesale price dislocation.

Sixth, the evaluation horizon for 2026 covers January--May 2026 (102 official OMIP trading sessions [decision epochs] and 46 non-trading dates [3 TARGET closures and 43 weekend days] across 148 calendar delivery days), corresponding to the data available at the time of study. Whether the performance advantages of \texttt{Dataset\_\allowbreak M8\_\allowbreak Physical\_\allowbreak NLP} over \texttt{Dataset\_\allowbreak NLP\_\allowbreak Subgroups} persist through subsequent seasonal cycles remains to be verified as additional market data become available.

Seventh, hyperparameter optimisation (Phase 4) was performed using a fixed seed ($s=42$) on the Hybrid representation and transferred to the Quantitative Baseline to isolate architectural capacity. While validated subsequently across multiple seeds, the HPO search itself reflects trajectory dynamics of that initialisation and naturally optimizes hyperparameters on the Hybrid feature space. Furthermore, the Phase 4 result is an exploratory post-selection capacity analysis over the 2026 horizon, not an unbiased out-of-sample evaluation. Finally, formal corrections for multiple comparisons (Bonferroni) were applied to the joint family of KS-tests but not to the multi-dataset ablation study in Phase 1, introducing the potential for family-wise error.

Eighth, while Section~\ref{subsec:trans_costs} establishes that net savings remain positive under proportional linear transaction fees up to $0.25$~\texteuro/MWh, the simulated environment does not model non-linear order book depth, market impact, or slippage during thin-volume trading sessions. Although all performance metrics are normalized to a 1~MW continuous baseload capacity ($E_{\text{load}, d} = 24$~MWh/day) and scale linearly, deploying large physical volumes (such as a 200~MW enterprise demand) on short-dated contracts ($D_1$--$D_3$) would be constrained by daily exchange liquidity and order-book depth. Future operational deployments should integrate execution cost models with live order-book depth to evaluate net procurement savings under strict market microstructure constraints.


